%----------------------------------------------------------------------------
\clearpage\section{Writing}\label{app:writing}
%----------------------------------------------------------------------------
\subsection{Plan before you write: problem to skeleton}\label{app:writing:spine}
Write the argument before the prose. We plan papers in this order, and each step builds on the previous one: problem; context and literature review; gap analysis; research questions; contributions; threat analysis; leading figure; examples; validation plan; spine; figures; skeleton; beats.
\begin{description}
  \item[Spine] One page that fixes the claims, the terminology and the notation of the whole paper. Every chapter, figure and review is checked against it. Its worst enemy is drift: terms, symbols and claims quietly changing between chapters.
  \item[Skeleton mind map] The argument made explicit before any prose: each chapter broken into paragraphs, each paragraph reduced to its topic sentence and a pointer to the evidence that backs it. Drafting then means filling in sentences around a skeleton already known to be sound. A mind map, on paper or in a tool, is the quickest way to draw it.
  \item[Beats] The moves inside each paragraph of the skeleton: what it establishes, which evidence it uses, and how it leads to the next one.
\end{description}
\advisortodo{Drafted from the paper-planning workflow, as the starter-pack note on the spine and skeleton mind map was empty; confirm, or add an example mind map.}

\subsection{The leading figure (vezérábra)}\label{app:writing:figure}
The leading figure (\cref{fig:overview}) explains the central idea in one picture: the problem, the contributions, and the result. Decide what a reader should take away from it, draw that, and write the caption so that the figure stands without the text. Draft it early, in the research phase, and treat it as exploratory until the spine is agreed: many papers find their argument while drawing it.
\advisortodo{Drafted from the paper-planning workflow, as the starter-pack note on the vezérábra was empty; confirm, or add good examples from past papers.}

\subsection{The abstract}
Follow \href{https://www.nature.com/documents/nature-summary-paragraph.pdf}{Nature's summary paragraph guide}; the abstract in this template already has its seven steps as placeholders.

\subsection{Style}
\begin{itemize}
  \item The Chicago School of Writing lectures: \href{https://www.youtube.com/watch?v=vtIzMaLkCaM}{part I}, \href{https://www.youtube.com/watch?v=aFwVf5a3pZM}{part II}, \href{https://www.youtube.com/watch?v=gphdfSwYMWU}{part III}.
  \item \href{https://www.youtube.com/watch?v=Unzc731iCUY}{MIT: How to Speak}, for the conference talk.
  \item Good examples: Leslie Lamport's papers and the Raft paper.
  \item Counterexamples: the Paxos and PBFT papers. Read them to see how hard a good idea is to follow when the writing does not help.
\end{itemize}

\subsection{Disclose AI use}
Disclose the use of AI tools in the form the TDK and thesis rules require; \cref{app:code:ai} covers the tools themselves.
